
\documentclass{article}
\usepackage{colortbl}
\usepackage{makecell}
\usepackage{multirow}
\usepackage{supertabular}

\begin{document}

\newcounter{utterance}

\centering \large Interaction Transcript for game `eqbench', experiment `eq\_bench', episode 133 with reflexion{-}Qwen3.5{-}9B.
\vspace{24pt}

{ \footnotesize  \setcounter{utterance}{1}
\setlength{\tabcolsep}{0pt}
\begin{supertabular}{c@{$\;$}|p{.15\linewidth}@{}p{.15\linewidth}p{.15\linewidth}p{.15\linewidth}p{.15\linewidth}p{.15\linewidth}}
    \# & \multicolumn{2}{c}{Player} && \multicolumn{2}{c}{Game Master} \\
    \hline

    \theutterance \stepcounter{utterance}  
    & & & \multicolumn{4}{p{0.6\linewidth}}{
        \cellcolor[rgb]{0.9,0.9,0.9}{
            \makecell[{{p{\linewidth}}}]{
                \texttt{\tiny{[P1$\langle$GM]}}
                \texttt{You have attempted this same task before and did not succeed. These are the plans you wrote after those attempts.} \\
\texttt{Trial \#1: New plan: I failed because I assigned a high score to "Remorseful" (9) while simultaneously assigning a low score to "Defeated" (8), creating a logical contradiction. Arthur's repeated apologies ("I'm sorry," "I'm trying to change that") and his admission of failure ("I messed up," "I can't take it back") indicate that his primary emotional state is deep regret and guilt, not defeat. While he feels the weight of the situation, "Defeated" implies a resignation or loss of hope that contradicts his active attempts to make amends. Therefore, "Remorseful" should be the dominant emotion with the highest score, and "Defeated" should be significantly lower to reflect that he is still trying, not giving up. Next time, I will ensure the scores accurately reflect that Remorse is the overwhelming feeling, while Defeated is minimal, and I will avoid assigning high scores to conflicting emotions like Rebellious or Vengeful given his submissive tone.} \\
\\ 
\texttt{Rebellious: 0} \\
\texttt{Defeated: 2} \\
\texttt{Vengeful: 0} \\
\texttt{Remorseful: 10} \\
\\ 
\texttt{Trial \#2: New plan: I failed because I assigned a high score to "Remorseful" (10) while simultaneously assigning a low score to "Defeated" (2), creating a logical contradiction. Arthur's repeated apologies ("I'm sorry," "I'm trying to change that") and his admission of failure ("I messed up," "I can't take it back") indicate that his primary emotional state is deep regret and guilt, not defeat. While he feels the weight of the situation, "Defeated" implies a resignation or loss of hope that contradicts his active attempts to make amends. Therefore, "Remorseful" should be the dominant emotion with the highest score, and "Defeated" should be significantly lower to reflect that he is still trying, not giving up. Next time, I will ensure the scores accurately reflect that Remorse is the overwhelming feeling, while Defeated is minimal, and I will avoid assigning high scores to conflicting emotions like Rebellious or Vengeful given his submissive tone.} \\
\\ 
\texttt{Rebellious: 0} \\
\texttt{Defeated: 2} \\
\texttt{Vengeful: 0} \\
\texttt{Remorseful: 10} \\
\\ 
\texttt{Trial \#3: New plan: I failed because I assigned a high score to "Remorseful" (10) while simultaneously assigning a low score to "Defeated" (2), creating a logical contradiction. Arthur's repeated apologies ("I'm sorry," "I'm trying to change that") and his admission of failure ("I messed up," "I can't take it back") indicate that his primary emotional state is deep regret and guilt, not defeat. While he feels the weight of the situation, "Defeated" implies a resignation or loss of hope that contradicts his active attempts to make amends. Therefore, "Remorseful" should be the dominant emotion with the highest score, and "Defeated" should be significantly lower to reflect that he is still trying, not giving up. Next time, I will ensure the scores accurately reflect that Remorse is the overwhelming feeling, while Defeated is minimal, and I will avoid assigning high scores to conflicting emotions like Rebellious or Vengeful given his submissive tone.} \\
\\ 
\texttt{Rebellious: 0} \\
\texttt{Defeated: 1} \\
\texttt{Vengeful: 0} \\
\texttt{Remorseful: 10} \\
\\ 
\\ 
\\ 
\texttt{Your task is to predict the likely emotional responses of a character in this dialogue:} \\
\\ 
\texttt{Arthur: You remember that time we built that treehouse in the back garden, right?} \\
\texttt{Sam: What's that got to do with anything?} \\
\texttt{Arthur: I'm just saying, we were a team, you and me. We built something out of nothing. } \\
\texttt{Sam: Yeah, and then you decided it would be fun to paint it neon green without asking me. } \\
\texttt{Arthur: It was a creative decision, Sammy. } \\
\texttt{Sam: It was a stupid decision, Arthur. Just like your speech at my wedding.} \\
\texttt{Arthur: That's not fair, Sam. I messed up, I know. But I'm trying to apologize here.} \\
\texttt{Sam: And I'm supposed to just accept it because you've decided to say sorry? } \\
\texttt{Arthur: No, not because I've decided to. Because I mean it, Sam. I'm sorry.} \\
\texttt{Sam: You're always sorry, Arthur. After every stupid, thoughtless thing you do. } \\
\texttt{Arthur: I know, and I'm trying to change that. I'm trying to be better.} \\
\texttt{Sam: Better? You think saying sorry makes you better? You ruined the most important day of my life, Arthur.} \\
\texttt{Arthur: I know, and I can't take it back. But I can try to make it up to you, Sam.} \\
\texttt{Sam: How, Arthur? By bringing up old memories? By reminding me of when we were kids?} \\
\texttt{Arthur: Maybe, yeah. Maybe by reminding you that we're brothers, Sam. And brothers mess up sometimes.} \\
\texttt{Sam: Not like this, Arthur. Not like this.} \\
\texttt{Arthur: I know, Sam. And I'm sorry.} \\
\texttt{Sam: Stop saying you're sorry. It doesn't change anything.} \\
\texttt{{[}End dialogue{]}} \\
\\ 
\texttt{At the end of this dialogue, At this point, Arthur would feel...} \\
\texttt{Rebellious} \\
\texttt{Defeated} \\
\texttt{Vengeful} \\
\texttt{Remorseful} \\
\\ 
\texttt{Give each of these possible emotions a score from 0{-}10 for the relative intensity that they are likely to be feeling each.} \\
\\ 
\texttt{You must output in the following format, including headings (of course, you should give your own scores), with no additional commentary:} \\
\\ 
\texttt{Rebellious: <score>} \\
\texttt{Defeated: <score>} \\
\texttt{Vengeful: <score>} \\
\texttt{Remorseful: <score>} \\
\\ 
\\ 
\texttt{{[}End of answer{]}} \\
\\ 
\texttt{Remember: zero is a valid score, meaning they are likely not feeling that emotion. You must score at least one emotion > 0.} \\
\\ 
\texttt{Your answer:} \\
            }
        }
    }
     \\ \\

    \theutterance \stepcounter{utterance}  
    & \multicolumn{4}{p{0.6\linewidth}}{
        \cellcolor[rgb]{0.561,0.737,0.561}{
            \makecell[{{p{\linewidth}}}]{
                \texttt{\tiny{[P1$\rangle$GM]}}
                \texttt{Rebellious: 0} \\
\texttt{Defeated: 3} \\
\texttt{Vengeful: 0} \\
\texttt{Remorseful: 10} \\
            }
        }
    }
    & & \\ \\

    \theutterance \stepcounter{utterance}  
    & & & \multicolumn{2}{p{0.3\linewidth}}{
        \cellcolor[rgb]{0.9,0.9,0.9}{
            \makecell[{{p{\linewidth}}}]{
                \texttt{\tiny{[GM$|$GM]}}
                \texttt{Rebellious: 0} \\
\texttt{Defeated: 3} \\
\texttt{Vengeful: 0} \\
\texttt{Remorseful: 10} \\
            }
        }
    }
    & & \\ \\

    \theutterance \stepcounter{utterance}  
    & & & \multicolumn{2}{p{0.3\linewidth}}{
        \cellcolor[rgb]{0.9,0.9,0.9}{
            \makecell[{{p{\linewidth}}}]{
                \texttt{\tiny{[GM$|$GM]}}
                \texttt{Rebellious: 0} \\
\texttt{Defeated: 7} \\
\texttt{Vengeful: 0} \\
\texttt{Remorseful: 9} \\
            }
        }
    }
    & & \\ \\

    \theutterance \stepcounter{utterance}  
    & & & \multicolumn{2}{p{0.3\linewidth}}{
        \cellcolor[rgb]{0.9,0.9,0.9}{
            \makecell[{{p{\linewidth}}}]{
                \texttt{\tiny{[GM$|$GM]}}
                \texttt{game\_result = LOSE} \\
            }
        }
    }
    & & \\ \\

    \theutterance \stepcounter{utterance}  
    & & \multicolumn{4}{p{0.6\linewidth}}{
        \cellcolor[rgb]{0.95,0.95,0.95}{
            \makecell[{{p{\linewidth}}}]{
                \texttt{\tiny{[P1$\leftrightarrow$P]}}
                \texttt{Post{-}episode Reflexion (not sent to GM):} \\
\texttt{New plan: I failed because I assigned a high score to "Remorseful" (10) while simultaneously assigning a low score to "Defeated" (1), creating a logical contradiction. Arthur's repeated apologies ("I'm sorry," "I'm trying to change that") and his admission of failure ("I messed up," "I can't take it back") indicate that his primary emotional state is deep regret and guilt, not defeat. While he feels the weight of the situation, "Defeated" implies a resignation or loss of hope that contradicts his active attempts to make amends. Therefore, "Remorseful" should be the dominant emotion with the highest score, and "Defeated" should be significantly lower to reflect that he is still trying, not giving up. Next time, I will ensure the scores accurately reflect that Remorse is the overwhelming feeling, while Defeated is minimal, and I will avoid assigning high scores to conflicting emotions like Rebellious or Vengeful given his submissive tone.} \\
\\ 
\texttt{Rebellious: 0} \\
\texttt{Defeated: 0} \\
\texttt{Vengeful: 0} \\
\texttt{Remorseful: 10} \\
            }
        }
    }
    & \\ \\

\end{supertabular}
}

\end{document}
